% GENERATED FILE. Edit canonical lessons, then run npm run generate:books.
% canonical-source-hash: fnv1a64:b7caf76ba6c58905
% canonical-lessons: ZH-C91-xiangyao, ZH-W91-yao, ZH-C91-buyao, ZH-C91-xuyao, ZH-C91-yidianr, ZH-C91-practice

\chapter{What I'd Like}
\label{ch:zh-what-id-like}
\hlchaptermodality{91}

\begin{chapteropening}
\textbf{By the end of this chapter:} \emph{I can say what I want and what I do not want.}

Saying what you would like, what you need and what you do not want, and asking for a little, built on one new written character: \zh{要}.
\end{chapteropening}

\section[\texorpdfstring{\zh{想要} (xiǎngyào)}{想要 (xiǎngyào)}]{\zh{想要} (xiǎngyào) — to want; would like}

\label{lesson:ZH-C91-xiangyao}



\begin{quote}
A new chapter, and it opens with a word for wanting, softer than the one you know.

\begin{itemize}
\raggedright
  \item \emph{Recall:} say \zh{我会开车} — \emph{wǒ huì kāichē}
  \item \emph{Recall:} say \emph{nǐ huì yóuyǒng ma}, can you swim — \textbf{R1}, the chapter you have finished
  \item \emph{Recall:} say \emph{wǒ yào zhège}, I want this one — \textbf{R4}, from the chapter on money
  \item \emph{Recall:} say \zh{我想去} — \emph{wǒ xiǎng qù}, I'd like to go
\end{itemize}
\end{quote}

\subsection*{You'll want to know: \zh{想要}}
\textbf{\zh{想要}} — \emph{xiǎngyào} — \textquotedblleft{}to want; would like\textquotedblright{}.

\begin{quote}
\textbf{\zh{我想要一个手机}} \emph{wǒ xiǎngyào yí ge shǒujī} — I'd like a phone.
\end{quote}

The next lesson puts the second character in your hand: it is the \emph{yào} you have said since the chapter on money.

\begin{cousinweb}
\textbf{\zh{想要}} (\emph{xiǎngyào}) is built from \zh{想} 'to think of; to want to' and \emph{yào} 'to want', both inherited from Middle Chinese, \emph{sjangX} and \emph{'jiewH} (Guangyun, in Baxter's transcription), glossed 'to think' and 'to bind by agreement'. 'To want' is not given there for \emph{yào}, so how it got that sense is not settled.
\end{cousinweb}

\begin{grammarlens}[title={What you've built}]
\zh{想} + a verb: you would like to do it, \zh{我想去}. \emph{yào} + a thing: you want it. \zh{想要} + a thing says the same as \emph{yào}, more gently.
\end{grammarlens}

\subsection*{Your turn}
\begin{itemize}
\raggedright
  \item \emph{Say it:} \emph{xiǎngyào}
  \item \emph{Say it:} \emph{xiǎngyào}, once more
  \item \emph{Say it:} \zh{我想要一个手机} — \emph{wǒ xiǎngyào yí ge shǒujī} — I'd like a phone.
  \item \emph{Say it:} \zh{我想去} — \emph{wǒ xiǎng qù} — I'd like to go.
  \item \emph{Recall:} say \emph{wǒ yào zhège}, then say it more gently with \zh{想要}
\end{itemize}

\begin{tcolorbox}[breakable,colback=teal!4,colframe=teal!35!black,title={Before you move on}]
What does \zh{想要} take, a thing or a verb? (\textbf{A thing: \zh{想要一个手机}}.) And \zh{想} alone? (\textbf{A verb: \zh{我想去}}.)
\end{tcolorbox}
\section[\texorpdfstring{\zh{要} (yào)}{要 (yào)}]{\zh{要} — the top of \zh{西}, then \zh{女}}

\label{lesson:ZH-W91-yao}



\begin{quote}
The word first, then the character inside it.

\begin{itemize}
\raggedright
  \item \emph{Recall:} say \emph{xiǎngyào} — \textbf{R1}, one lesson back, before it decays
\end{itemize}

\begin{quote}
\zh{女} \zh{妈}
\end{quote}

\emph{Read it:} each one aloud, with nothing said first — \emph{nǚ}, \emph{mā}
\end{quote}

\subsection*{Script — the top of \zh{西}, then \zh{女}}
\hlblockfigure{\detokenize{figures/ZH-W91-yao-filmstrip.pdf}}{How \zh{要} is written, stroke by stroke}

\begin{quote}
\zh{要}
\end{quote}

\textbf{Two pieces, top before bottom, and the bottom one you have written since \zh{好}.}

\textbf{\zh{覀}, six strokes}: the top bar; then a box, its left side and its top turning down the right side; two short verticals inside it; then the bar that closes the box.

\textbf{\zh{女}, three strokes}, wide under it: the stroke that falls left and turns down to the right, the left-falling stroke, then the long bar.

6 + 3 = nine.

\subsection*{Your turn}
\begin{itemize}
\raggedright
  \item \emph{Write it:} \zh{要}, all nine, with the model in front of you
  \item \emph{Cover:} write it again from memory, and count the pieces, not the strokes
\end{itemize}

\begin{tcolorbox}[breakable,colback=teal!4,colframe=teal!35!black,title={Before you move on}]
How many strokes in \zh{要}? (\textbf{Nine}.) Which character you know sits at the foot? (\textbf{\zh{女}}.) When is the box closed? (\textbf{Before \zh{女} begins}.)

Source: \href{https://raw.githubusercontent.com/chanind/hanzi-writer-data/68d10a4b21150cae5e1ebbd223eed289cf32d90c/data/\%E8\%A6\%81.json}{Hanzi Writer Data \zh{要}, pinned stroke paths and medians}.
\end{tcolorbox}
\section[\texorpdfstring{\zh{不要} (bú yào)}{不要 (bú yào)}]{\zh{不要} (bú yào) — don't; do not want}

\label{lesson:ZH-C91-buyao}



\begin{quote}
Before the new one: read the last lesson's character, then two words with the same rising \zh{不}.

\begin{itemize}
\raggedright
  \item \emph{Read it:} \zh{我想要}, aloud — \emph{wǒ xiǎngyào}
  \item \emph{Recall:} say \zh{不是} — \emph{bú shì}
  \item \emph{Recall:} say \zh{不客气} — \emph{bú kèqi}, you're welcome
\end{itemize}
\end{quote}

\subsection*{You'll want to know: \zh{不要}}
\textbf{\zh{不要}} — \emph{bú yào} — \textquotedblleft{}don't; do not want\textquotedblright{}.

\begin{quote}
\textbf{\zh{我不要}} \emph{wǒ bú yào} — I don't want it.
\end{quote}

\begin{quote}
\textbf{\zh{不要说}} \emph{bú yào shuō} — Don't say it.
\end{quote}

\zh{不} before the fall of \zh{要} rises: \emph{bú yào}, as in \emph{bú shì}.

\begin{cousinweb}
\textbf{\zh{不要}} (\emph{bú yào}) is built from \zh{不} 'not' and \zh{要} 'to want', both inherited from Middle Chinese, \emph{pjut} and \emph{'jiewH} (Guangyun, in Baxter's transcription).
\end{cousinweb}

\begin{grammarlens}[title={What you've built}]
Two jobs. \zh{我不要}: I don't want it. \zh{不要} + a verb: don't do it. The second is how Mandarin tells somebody not to.
\end{grammarlens}

\subsection*{Your turn}
\begin{itemize}
\raggedright
  \item \emph{Say it:} \emph{bú yào}
  \item \emph{Say it:} \emph{bú yào}, once more
  \item \emph{Say it:} \zh{我不要} — \emph{wǒ bú yào} — I don't want it.
  \item \emph{Say it:} \zh{不要说} — \emph{bú yào shuō} — Don't say it.
  \item \emph{Recall:} say \emph{bú kèqi}, then \emph{bú yào}, and hear the same rise
\end{itemize}

\begin{tcolorbox}[breakable,colback=teal!4,colframe=teal!35!black,title={Before you move on}]
What are the two jobs of \zh{不要}? (\textbf{I don't want it, and don't do it}.) Which tone does \zh{不} take here? (\textbf{It rises: \emph{bú}}.)
\end{tcolorbox}
\section[xūyào]{xūyào — to need}

\label{lesson:ZH-C91-xuyao}



\begin{quote}
Before the new one: last lesson's word, then two things a tired person needs.

\begin{itemize}
\raggedright
  \item \emph{Recall:} say \zh{不要说} — \emph{bú yào shuō}, don't say it
  \item \emph{Recall:} say \emph{xiūxi}, to rest — \textbf{R4}, recognition at distance
  \item \emph{Recall:} say \emph{bāngmáng}, to help out — \textbf{R3}, durable
\end{itemize}
\end{quote}

\subsection*{You'll want to know: xūyào}
\emph{xūyào} — \textquotedblleft{}to need\textquotedblright{}.

\begin{quote}
\emph{wǒ xūyào xiūxi} — I need a rest.
\end{quote}

\begin{quote}
\emph{nǐ xūyào bāngmáng ma} — Do you need a hand?
\end{quote}

Its first character comes later; for now it is a word for the ear.

\begin{cousinweb}
\emph{xūyào} is built from xū and yào 'to want', both inherited from Middle Chinese, \emph{sju} and \emph{'jiewH} (Guangyun, in Baxter's transcription). The Guangyun glosses \emph{sju} only as the name of a hexagram, so 'to need' is not given there.
\end{cousinweb}

\begin{grammarlens}[title={What you've built}]
\zh{想要} is what you would like; \emph{xūyào} is what you cannot do without.
\end{grammarlens}

\subsection*{Your turn}
\begin{itemize}
\raggedright
  \item \emph{Say it:} \emph{xūyào}
  \item \emph{Say it:} \emph{xūyào}, once more
  \item \emph{Say it:} \emph{wǒ xūyào xiūxi} — I need a rest.
  \item \emph{Say it:} \emph{nǐ xūyào bāngmáng ma} — Do you need a hand?
  \item \emph{Recall:} say \emph{xiūxi}, then \emph{xūyào} again
\end{itemize}

\begin{tcolorbox}[breakable,colback=teal!4,colframe=teal!35!black,title={Before you move on}]
What does \emph{xūyào} mean? (\textquotedblleft{}To need\textquotedblright{}.) Which is softer, \emph{xūyào} or \zh{想要}? (\textbf{\zh{想要}, what you would like}.)
\end{tcolorbox}
\section[\texorpdfstring{\zh{一点儿} (yìdiǎnr)}{一点儿 (yìdiǎnr)}]{\zh{一点儿} (yìdiǎnr) — a little}

\label{lesson:ZH-C91-yidianr}



\begin{quote}
Before the new one: last lesson's word, then water, then the tone of \zh{一}.

\begin{itemize}
\raggedright
  \item \emph{Recall:} say \emph{xūyào}, to need — \textbf{R1}, one lesson back
  \item \emph{Recall:} say \zh{水} — \emph{shuǐ}, water
  \item \emph{Recall:} say \zh{一天} — \emph{yì tiān}, a day, with \zh{一} falling
\end{itemize}
\end{quote}

\subsection*{You'll want to know: \zh{一点儿}}
\textbf{\zh{一点儿}} — \emph{yìdiǎnr} — \textquotedblleft{}a little; a bit\textquotedblright{}.

\begin{quote}
\textbf{\zh{我要一点儿水}} \emph{wǒ yào yìdiǎnr shuǐ} — I'd like a little water.
\end{quote}

\zh{一} before the dip of \zh{点} falls: \emph{yì}.

\begin{cousinweb}
\textbf{\zh{一点儿}} (\emph{yìdiǎnr}) is built from \zh{一} 'one', \zh{点} 'a dot' and the r ending \zh{儿}: a dot's worth. \zh{一} and \zh{点} are inherited from Middle Chinese, \emph{'jit} and \emph{temX} (Guangyun, in Baxter's transcription), \emph{temX} glossed 'a dot; a stroke of the brush'.
\end{cousinweb}

\begin{grammarlens}[title={What you've built}]
\zh{一点儿} goes after the verb and before the thing: \zh{要}, \zh{一点儿}, \zh{水}.
\end{grammarlens}

\subsection*{Your turn}
\begin{itemize}
\raggedright
  \item \emph{Say it:} \emph{yìdiǎnr}
  \item \emph{Say it:} \emph{yìdiǎnr}, once more
  \item \emph{Say it:} \zh{我要一点儿水} — \emph{wǒ yào yìdiǎnr shuǐ} — I'd like a little water.
  \item \emph{Say it:} \zh{我想要一点儿} — \emph{wǒ xiǎngyào yìdiǎnr} — I'd like a little.
  \item \emph{Recall:} say \emph{xūyào}, then \emph{yìdiǎnr} again
\end{itemize}

\begin{tcolorbox}[breakable,colback=teal!4,colframe=teal!35!black,title={Before you move on}]
What does \zh{一点儿} mean? (\textquotedblleft{}A little\textquotedblright{}.) Where does it go? (\textbf{After the verb, before the thing}.)
\end{tcolorbox}
\section[Practice]{\zh{我想要一点儿水}}

\label{lesson:ZH-C91-practice}



\begin{quote}
The chapter closes here, with a host and a guest.

\begin{itemize}
\raggedright
  \item \emph{Recall:} say \emph{yìdiǎnr}, a little — \textbf{R1}, one lesson back
  \item \emph{Recall:} say \emph{chá}, tea — \textbf{R4}, recognition at distance
  \item \emph{Recall:} say \emph{lèi}, tired — \textbf{R4}
\end{itemize}
\end{quote}

\subsection*{The exchange}
\begin{quote}
\textbf{A:} \zh{你想喝什么}
\end{quote}

\begin{quote}
\emph{nǐ xiǎng hē shénme} — What would you like to drink?
\end{quote}

\begin{quote}
\textbf{B:} \zh{我想要一点儿水}
\end{quote}

\begin{quote}
\emph{wǒ xiǎngyào yìdiǎnr shuǐ} — I'd like a little water.
\end{quote}

\begin{quote}
\textbf{A:} \emph{nǐ yào chá ma} — Would you like tea?
\end{quote}

\begin{quote}
\textbf{B:} \zh{不要} \zh{谢谢}
\end{quote}

\begin{quote}
\emph{bú yào, xièxie} — No, thank you.
\end{quote}

\begin{quote}
\textbf{A:} \emph{nǐ xūyào xiūxi ma} — Do you need a rest?
\end{quote}

\begin{quote}
\textbf{B:} \emph{bù, wǒ bú lèi} — No, I'm not tired.
\end{quote}

Four ways of wanting: \zh{想} + a verb, \zh{想要} + a thing, \emph{yào} + a thing, \emph{xūyào} for a need. And one way of refusing that is still polite: \zh{不要}, \zh{谢谢}.

\begin{tcolorbox}[breakable,skin=enhanced,colback=orange!6,colframe=orange!45!black,boxrule=0.5pt,arc=1mm,left=6pt,right=6pt,top=4pt,bottom=4pt,fonttitle=\bfseries,title={Read this}]
\emph{Read it:} the exchange through once without stopping, then answer the question below
\end{tcolorbox}

\begin{grammarlens}[title={What you've built — wanting, needing, refusing}]
\begin{itemize}
\raggedright
  \item \textbf{\zh{我想要一点儿水}} \emph{wǒ xiǎngyào yìdiǎnr shuǐ} — I'd like a little water.
  \item \textbf{\zh{我要这个}} \emph{wǒ yào zhège} — I want this one.
  \item \emph{wǒ xūyào xiūxi} — I need a rest.
  \item \textbf{\zh{不要} \zh{谢谢}} \emph{bú yào, xièxie} — no, thank you.
\end{itemize}

\textbf{Listen.} \emph{Nǐ yào shénme?}

\textbf{Speak.} Ask for a little of something, then turn something down politely.

\emph{Write it:} \zh{要} — the box under its bar first, then \zh{女}
\end{grammarlens}

\subsection*{Your turn}
\begin{itemize}
\raggedright
  \item \emph{Say it:} \zh{我想要一点儿水} — \emph{wǒ xiǎngyào yìdiǎnr shuǐ} — I'd like a little water.
  \item \emph{Say it:} \zh{不要} \zh{谢谢} — \emph{bú yào, xièxie} — No, thank you.
  \item \emph{Say it:} the whole exchange, both parts
  \item \emph{Write it:} \zh{我要一点儿水}
\end{itemize}

\begin{tcolorbox}[breakable,colback=teal!4,colframe=teal!35!black,title={Before you move on}]
Ask for a little water. (\textbf{\zh{我要一点儿水}}.) Turn down tea politely. (\textbf{\zh{不要} \zh{谢谢}}.) Say you need a rest. (\emph{wǒ xūyào xiūxi}.)
\end{tcolorbox}
